\section{Delovanje mehurčne komore}

Izumitelj mehurčne komore, D. A. Glaser, je predpostavil, da je za nastanek mehurčkov kriva elektrostatska sila med električnimi naboji, ki jih povzroči ioniziran delec\cite{glaser_bubble_1958}. Glaser je sčasoma opustil svojo teorijo\cite{seitz_theory_1958}. V seminarju bom predstavil dve od možnih teorij, kako nastanejo mehurčki v mehurčni komori. Zaradi narave uporabe mehurčne komore, tj. visokoenergijske raziskave, fokus nikoli ni bil na njenem delovanju, ampak na možnih novih spoznanjih kot posledica raziskav\cite{seitz_theory_1958}. Prva teorija predstavi termalni model, ki ga je oblikoval Seitz, in predpostavlja nastanek ``termalni skokov'' (ang. \textit{thermal spikes})\cite{seitz_theory_1958}. Druga teorija je novejša, iz leta 1998, in nastanek mehurčkov v vodikovih in helijevih komorah pripisuje mehanskemu procesu\cite{konstantinov_how_1998}. Seitzov model je trenutno sprejet kot najboljša razlaga nukleacije mehurčkov kot posledice sevanja v pregretih tekočinah\cite{bruno_critical_2019}.

\subsection{Termalni model}

Komora je napolnjena s pregreto kapljevino kot so helij, devterij, vodik, propan, freon ali ksenon\footnote{Glaser je v komunikaciji s Seitzom predstavil dokaze, da zelo čist(?) ksenon ni občutljiv na sevanje, vendar postane ob dodajanju malih količin etilena.}. Spomnimo, da je pregreta kapljevina kapljevina, ki smo jo počasi segreli nad njeno vrelišče pri konstantnem tlaku, brez da bi kapljevina zavrela. To lahko storimo, če je stena posode gladka, kapljevina miruje in v njej ni drugih delcev\cite{strnad_fizika_2016}. Metastabilno stanje pregrete kapljevine se popravi v stabilno v trenutku, ko nastane v njej mehurček, kapljevina pa eksplozivno vzkipi. Temperatura kapljevine pade na vrelišče.



\subsection{Mehanski model}

